\documentclass[a4paper,12pt]{article}
\usepackage[utf8]{vietnam}
\usepackage{amsmath, amssymb, graphicx}
\usepackage{listings}
\usepackage{xcolor}
\usepackage{fancyhdr}
\usepackage{url}
\usepackage{booktabs}

% Cấu hình header & footer
\pagestyle{fancy}
\fancyhf{}
\renewcommand{\headrulewidth}{0.5pt}
\renewcommand{\footrulewidth}{0.5pt}

% Cấu hình màu cho code
\definecolor{codegray}{gray}{0.9}
\lstset{
	backgroundcolor = \color{codegray},
	basicstyle = \ttfamily\footnotesize,
	breaklines = true,
	frame = single,
}

\begin{document}

\title{Ứng dụng Cơ chế AHRRT trong Lập kế hoạch Quỹ đạo trên Địa hình 2.5D}
\author{Đức Anh \\ Nhóm 2}
\maketitle

\setcounter{section}{0}

\section{Lời nói đầu: Sự chuyển đổi từ 3D sang 2.5D}
Thuật toán \textbf{AHRRT} (Adaptive RRT with Heuristic Search) \cite{202511.1805} ban đầu được thiết kế cho xe tự hành (2D) và UAV trong khối đô thị (3D). Tuy nhiên, khi ứng dụng vào \textbf{địa hình 2.5D} (mặt phẳng địa hình đồi núi, tòa nhà nới màn hình biểu diễn dưới dạng bản đồ độ cao Heightmap $z = f(x,y)$), cơ chế của AHRRT phải được tùy biến lại. Dưới đây là rãnh phân tích về cách hoạt động, ảnh hưởng và kỹ thuật triển khai AHRRT trong bối cảnh này.

%---------------------------------------------------------
\section{Cách hoạt động của AHRRT trên nền 2.5D}

Bốn chiến lược cốt lõi của AHRRT khi áp đặt lên một bề mặt địa hình 2.5D (bề mặt nhấp nhô) hoạt động như sau:

\subsection{1. Kích thước bước nhảy thích ứng theo "Độ dốc" (Adaptive Step)}
Ở không gian 3D, bước nhảy co thụt theo khoảng cách rỗng tĩnh tới vật cản. 
Trong 2.5D, \textbf{vật cản chính là yếu tố cao độ}. Khi thuật toán mở rộng nhánh lên một vùng có độ dốc (gradient cao) gồ ghề, kích thước bước nhảy $tem\_step$ sẽ tự động tính toán co ngắn lại để luồn lách theo khe núi. Ngược lại, robot tiến ra vùng bình nguyên (độ cao $Z$ dao động ít) thì bước nhảy bung rộng tối đa để đi lướt thật nhanh.

\subsection{2. Thiên hướng Mục tiêu có trọng số cao độ (Target Bias)}
Thông thường quy tắc sẽ luôn kéo $\vec{X}_{new}$ thẳng về Đích $\vec{X}_{goal}$. 
Trong 2.5D, đường nhắm trên trục toạ độ mặt bằng (x, y) là thẳng đỉnh, nhưng thuật toán sẽ tự đánh tháo nếu vector đâm thẳng này yêu cầu robot phải "leo thẳng đứng (Z)" vượt qua một đỉnh núi tốn khối lượng lớn năng lượng pin của UAV.  

\subsection{3. Lực trường thế đẩy từ các đỉnh cao (APF Attraction-Repulsion)}
Trường thế (APF) trong 2.5D hoạt động vô cùng hiệu quả:
\begin{itemize}
    \item \textbf{Lực đẩy $\vec{F}_{rep}$}: Được phát ra mạnh mẽ từ các \textbf{đỉnh núi nhô cao} hoặc công trình che khuất tầm bay. Tâm của lực đẩy được chiếu trên ma trận lưới cao độ. Trọng lực sẽ đẩy văng hướng quy hoạch sang các sườn ngầm thung lũng (Z thấp).
    \item \textbf{Lực hút $\vec{F}_{att}$}: Chịu trách nhiệm kéo tụt agent bám vào đường hẻm dốc dẫn hướng thoát thung để về đến mặt bằng đích đến một cách bằng phẳng nhẹ đà nhất.
\end{itemize}

\subsection{4. Phẫu thuật cắt tỉa bám theo Line of Sight (LOS)}
Thuật toán AHRRT nguyên thủy nối thử 2 điểm trung gian rời rạc để kiểm tra có va chạm cản khối hay không. 
Ở bản đồ 2.5D, quá trình Cắt Tỉa (Pruning) chính là kiểm tra \textbf{Tầm nhìn thẳng (Line of Sight)}. Một đường bay được gọt thẳng đi loại bỏ mốc nếu toàn bộ phân đoạn quỹ đạo nối hai điểm trên không hoàn toàn \textit{cao hơn} mặt đất bên dưới:
\begin{equation}
Z_{path} > Heightmap(x_i, y_i) + Clearance
\end{equation}

%---------------------------------------------------------
\section{Ảnh hưởng và Hiệu năng của AHRRT trong 2.5D}

\begin{itemize}
    \item \textbf{Tiết kiệm Năng Lượng (Energy Efficiency):} Nhờ thao tác phân bố kết hợp APF đánh né xa cấu trúc chóp cản trở, quỹ đạo tránh được các động tác vòng xoáy gắt qua đó tối ưu lượng pin cho cỗ máy thao tác theo địa hình.
    \item \textbf{Khử Triệt Răng Cưa Đường Đi:} Drone hay robot thường bay leo núi gặp chứng lởm chởm vì thuật toán nguyên thủy nội suy lên xuống theo từng mắt lưới. AHRRT nhổ sạch gai xương nhờ có chức năng làm trơn đồ thị bằng khâu cắt điểm gãy khúc.
    \item \textbf{Tốc Độ Hội Tụ Thần Tốc:} Lực hút Target Bias \cite{202511.1805} hỗ trợ xua đuổi thuật toán không bị loang vũng (lan man tính toán cả ở phần đuôi phía sau), giúp UAV trong môi trường rủi ro chập trùng kịp tính lối tẩu thoát ở môi trường 2.5D phức hợp giới hạn băng thông.
\end{itemize}

%---------------------------------------------------------
\section{Cách sử dụng và cấu hình để triển khai trong 2.5D}

Để mô đun hoạt động trên Heightmap (2.5D), ta cần làm những bước sau trước khi bấm Start vào lõi AHRRT:

\subsection{1. Thay đổi hàm đo lường chi phí (Cost Function)}
Khoảng cách trong 2.5D không nên đo mù quáng theo $Euclidean$ không gian, vì việc thay đổi cao độ cắn dứt năng lượng bay hơn việc bay tiến rất nhiều (khắc phục quy luật leo đồi). Cần điều chế trong hàm tìm kiếm Nearest Neighbor công thức phạt $Z$:
\begin{equation}
Distance = w_1 \sqrt{(x_2 - x_1)^2 + (y_2 - y_1)^2} + w_2 |z_2 - z_1|
\end{equation}
Với $w_2 \gg w_1$ để nhắc khéo AHRRT ưu tiên rúc đường uốn lượn thung lũng (lực đẩy của APF).

\subsection{2. Điều chỉnh lại hàm CollisionCheck() theo chuẩn 2.5D}
Trái ngược 3D cần dò giao điểm (Object Intersection 3D Volume), 2.5D biến quá trình kiểm tra vật cản thành ma trận O(1) rất nhanh. Hàm \texttt{CollisionFree(point)} trong cấu hình này trở thành: \textit{"Kiểm tra xem Cao độ Z của thuật toán tại (x,y) đang quy hoạch có vượt quá chiều cao điểm lấy của mắt lưới cộng bù viền an toàn bay hay không."}

\subsection{3. Biến thiên Parameter p (Xác suất Target Bias)}
Trong núi nhấp nhô 2.5D, bị rơi vào mỏm cục bộ lõm thung lũng (local minima) là cực dễ. Tham số phân vùng ưu tiên $P_{target}$ của AHRRT nên chuyển thành kịch bản linh động: Ở mặt đồng thung, xả $P = 0.8$ nhằm chạy siêu lẹ. Bị kẹt trong góc cụt hẻm núi -> ép lại $P = 0.1$ để cây RRT lan rộng trổ mầm ra ngoài ngõ cụt vách đá trước rồi tính tiếp.

\bibliographystyle{IEEEtran}
\bibliography{refs}

\end{document}